\documentclass[UTF8,a4paper,12pt]{ctexart}

\usepackage[left=2.6cm,right=2.6cm,top=2.5cm,bottom=2.5cm]{geometry}
\usepackage{booktabs}
\usepackage{longtable}
\usepackage{array}
\usepackage{xcolor}
\usepackage{listings}
\usepackage{enumitem}
\usepackage{hyperref}
\usepackage{fancyhdr}
\usepackage{graphicx}

\hypersetup{colorlinks=true,linkcolor=black,urlcolor=blue}
\setlength{\parindent}{2em}
\setlength{\parskip}{0.25em}
\setlength{\headheight}{15pt}
\setlength{\emergencystretch}{3em}
\renewcommand{\arraystretch}{1.25}

\pagestyle{fancy}
\fancyhf{}
\fancyhead[L]{数据结构实习报告}
\fancyhead[R]{电梯离散事件模拟}
\fancyfoot[C]{\thepage}

\definecolor{codegreen}{RGB}{0,110,50}
\definecolor{codegray}{RGB}{90,90,90}
\definecolor{codeblue}{RGB}{20,70,160}
\definecolor{backcolor}{RGB}{247,247,247}

\lstdefinestyle{cstyle}{
    language=C,
    backgroundcolor=\color{backcolor},
    basicstyle=\ttfamily\small,
    keywordstyle=\color{codeblue}\bfseries,
    commentstyle=\color{codegreen},
    stringstyle=\color{red!60!black},
    numbers=left,
    numberstyle=\tiny\color{codegray},
    numbersep=8pt,
    breaklines=true,
    breakatwhitespace=false,
    showstringspaces=false,
    tabsize=4,
    frame=single,
    columns=fullflexible,
    keepspaces=true
}
\lstset{style=cstyle}

\title{\heiti 电梯离散事件模拟系统\\[0.5em]\Large 实习报告}
\author{[REDACTED]\\学号：[REDACTED]}
\date{\today}

\begin{document}

\maketitle
\thispagestyle{empty}
\vfill
\begin{center}
\begin{tabular}{rl}
实验内容：& 单电梯及双电梯离散事件模拟\\
语言：& C\\
数据结构：& 有序事件链表、链队列\\
时间单位：& $t=0.1\,\mathrm{s}$
\end{tabular}
\end{center}
\vfill
\clearpage

\tableofcontents
\clearpage

\section{需求分析}

\subsection{问题描述}

本项目模拟一栋五层教学楼中的电梯运行。楼层编号为0至4，其中1层是电梯空闲时的本垒层。乘客在给定时刻到达某层，选择目标楼层并按下上行或下行按钮；若等待超过个人可容忍时间，则放弃乘梯。电梯根据楼层请求和轿厢请求完成开门、人员进出、关门、加速、逐层运行、减速和方向切换等动作。

系统采用离散事件模拟方法。模拟时钟不按固定步长循环，而是每次从事件表中取出发生时间最早的事件，将系统时间直接推进到该时刻，再执行该事件并插入由它产生的后继事件。这样既能够保持动作的时序关系，也避免在两个事件之间进行无意义的空循环。

在完成单电梯模拟后，项目进一步实现两部电梯协同运行。双电梯共享楼层等待队列与厅外按钮，但各自维护所在楼层、运行方向、轿厢按钮和轿厢乘客。系统需要把每个厅外请求唯一分配给一部电梯，避免两部电梯重复响应同一请求。

\subsection{功能需求}

\begin{enumerate}[label=\arabic*.]
    \item 模拟0至4层的楼层环境，初始时电梯位于1层并处于空闲状态。
    \item 按发生时间处理乘客事件M1至M6及电梯事件E1至E9。
    \item 支持乘客到达、排队、等待超时、进入电梯和到达目标层离开等行为。
    \item 支持厅外上行、下行按钮及轿厢内部目标楼层按钮。
    \item 开门和关门各耗时20t，人员每次进出耗时25t；电梯每40t检查一次是否可以关门。
    \item 电梯上行一层耗时51t，下行一层耗时61t，并考虑加速、减速事件的时间开销。
    \item 电梯空闲超过300t且不在1层时返回1层待命。
    \item 输出所有事件的时间、事件类型、乘客编号、电梯编号、楼层及运行状态。
    \item 双电梯模式下，根据距离、方向与轿厢负载选择响应电梯。
    \item 程序结束时释放事件链表和全部队列占用的动态内存。
\end{enumerate}

\subsection{输入与输出}

测试采用15位固定乘客。每位乘客包含编号、到达时刻、进入楼层、目标楼层和最长等待时间。模拟终止时间设置为6000t，可以覆盖最后一位乘客在4000t到达后的完整处理过程。

\begin{longtable}{ccccc}
\toprule
乘客ID & 到达时间/t & 进入楼层 & 目标楼层 & 最长等待/t\\
\midrule
\endfirsthead
\toprule
乘客ID & 到达时间/t & 进入楼层 & 目标楼层 & 最长等待/t\\
\midrule
\endhead
1&0&1&4&500\\
2&100&0&3&300\\
3&150&2&0&500\\
4&200&3&1&500\\
5&250&4&0&500\\
6&300&1&2&100\\
7&400&0&4&500\\
8&500&2&3&500\\
9&1000&1&0&500\\
10&1500&3&2&500\\
11&2000&4&1&500\\
12&2500&0&4&200\\
13&3000&2&1&500\\
14&3500&1&3&500\\
15&4000&3&0&500\\
\bottomrule
\end{longtable}

输出以事件日志为主，例如：

\begin{lstlisting}[language={},numbers=none]
Time=741 Person=8 Left Elevator=1 at floor 3
Time=5209 Event=E1 Elevator=0 Floor=1 State=Idle
Simulation finished at Time=5209
Left=13 GaveUp=2 Total=15
\end{lstlisting}

\section{概要设计}

\subsection{总体结构}

系统分为链队列模块、事件表模块、乘客状态模块、电梯状态模块、调度模块和主程序模块。主程序只负责初始化测试数据和驱动事件循环；具体逻辑集中在操作函数中。

\begin{center}
\begin{tabular}{p{3.5cm}p{10cm}}
\toprule
模块 & 功能\\
\midrule
链队列 & 提供队列初始化、入队、出队、判空、求长度和销毁操作\\
事件表 & 以时间升序保存未来事件，支持稳定插入、弹出和指定事件删除\\
乘客 & 处理M1至M6，维护乘客生命周期与楼层等待队列\\
电梯 & 处理E1至E9，维护门状态、运行方向、楼层和轿厢请求\\
调度 & 为厅外请求选择电梯，维护请求与电梯之间的分配关系\\
主程序 & 初始化15位乘客，循环弹出事件，完成分派并统计结果\\
\bottomrule
\end{tabular}
\end{center}

\subsection{离散事件循环}

主程序按照以下步骤运行：

\begin{enumerate}[label=\arabic*.]
    \item 初始化模拟器、楼层等待队列、轿厢队列和事件表；
    \item 将15位乘客的M1事件按到达时刻插入事件表；
    \item 弹出事件表头结点，并把模拟时钟更新为该事件时刻；
    \item 根据事件编号把M1至M6交给\texttt{PersonState}，把E1至E9交给\texttt{ElevatorState}；
    \item 事件处理函数修改系统状态，并把新的未来事件插回事件表；
    \item 当事件表为空或事件时间超过6000t时结束模拟；
    \item 统计乘客结果并销毁所有动态结构。
\end{enumerate}

\subsection{单电梯与双电梯结构比较}

单电梯版本把\texttt{floor}、\texttt{state}、\texttt{CallCar}和\texttt{InElevator}直接放在\texttt{Simulator}中。双电梯版本增加\texttt{Elevator}结构，并由\texttt{Simulator}保存电梯数组：

\begin{lstlisting}
typedef struct Elevator{
    int No;
    int floor;
    int D1,D2,D3;
    int state;
    int CallCar[5];
    LinkQueue InElevator;
}Elevator;

typedef struct Simulator{
    int time;
    Elevator elevators[ELEVATOR_COUNT];
    LinkQueue waitingQueue[5];
    int CallUp[5],CallDown[5];
    int assignedUp[5],assignedDown[5];
    EventList events;
}Simulator;
\end{lstlisting}

其中，楼层等待队列与厅外按钮由所有电梯共享；轿厢按钮和轿厢乘客队列属于具体电梯。\texttt{assignedUp}和\texttt{assignedDown}记录每个厅外请求被分配给的电梯编号，$-1$表示尚未分配。

\section{详细设计}

\subsection{核心数据结构}

\subsubsection{事件结点}

事件结点保存事件发生时刻、类型、乘客编号、楼层和电梯编号。单电梯版本不需要\texttt{elevatorNo}；双电梯版本使用$-1$表示与具体电梯无关的乘客事件。

事件表采用单链表并按\texttt{time}升序排列。插入时使用小于等于比较，因此同一时刻的事件保持插入先后次序，保证模拟结果可重复。

\subsubsection{乘客状态}

双电梯版本把乘客生命周期划分为以下状态：

\begin{center}
\begin{tabular}{cl}
\toprule
状态 & 含义\\
\midrule
NotArrived & 尚未进入系统\\
Waiting & 已进入楼层等待队列\\
Assigned & 已安排进入某部电梯，但M5尚未执行\\
In\_Elevator & 已进入指定电梯\\
Left & 已到达目标楼层并离开\\
GaveUp & 等待超时并放弃\\
\bottomrule
\end{tabular}
\end{center}

增加\texttt{Assigned}状态可以处理“已经安排M5但尚未真正进入”的时间间隔，并使过期M4事件能够通过状态判断安全失效。

\subsubsection{队列}

系统复用已有的单链队列实现。五个\texttt{waitingQueue}分别保存各层等待乘客编号，每部电梯的\texttt{InElevator}保存轿厢内乘客编号。队列元素只保存整数编号，详细乘客信息通过\texttt{people[personNo]}访问。

已有\texttt{DeQueue}只能删除队首，而超时乘客可能位于队列中间，因此新增\texttt{personDelete}，通过前驱指针查找并删除指定编号，同时在删除队尾时修正\texttt{rear}指针。

\subsection{事件表操作}

\texttt{InsertEvent}按时间插入事件，时间复杂度为$O(n)$；\texttt{PopEvent}删除表头，复杂度为$O(1)$；\texttt{DeleteEvent}按事件类型和电梯编号删除过期事件，复杂度为$O(n)$。多电梯版本的事件判定必须同时比较\texttt{type}和\texttt{elevatorNo}，否则取消0号电梯的E7时会误删1号电梯的E7。

\subsection{乘客事件设计}

\begin{longtable}{cp{11cm}}
\toprule
事件 & 处理内容\\
\midrule
M1 & 乘客进入模拟系统，初始化其状态并立即安排M2\\
M2 & 乘客按下楼层按钮，随后进入排队处理M3\\
M3 & 加入所在层等待队列，设置上行或下行按钮，分配响应电梯，并安排超时事件M4\\
M4 & 若乘客仍为Waiting或Assigned，则从等待队列删除并标记GaveUp；若已进入电梯则忽略该过期事件\\
M5 & 从楼层队列删除乘客，加入指定电梯的轿厢队列，设置目标层CallCar\\
M6 & 从指定电梯删除乘客，刷新轿厢按钮并标记Left\\
\bottomrule
\end{longtable}

放弃事件的发生时刻为
\[
T_{M4}=T_{arrival}+T_{giveup}.
\]
M5和M6执行后通过\texttt{personEventTrue}判断当前电梯是否还有人员进出事件；如果没有，则令$D1=0,D3=1$，表示门仍开着但已无人进出。

\subsection{电梯事件设计}

\begin{longtable}{cp{11cm}}
\toprule
事件 & 处理内容\\
\midrule
E1 & 空闲等待；超过300t且不在1层时生成返回1层的轿厢请求\\
E2 & 到达一层后的减速判断；有本层请求则安排E3，否则重新调用控制器\\
E3 & 开门并设置门状态，安排E4、E5和300t后的E9\\
E4 & 按进入电梯的逆序安排离开事件，再按队列顺序安排进入事件\\
E5 & 每40t检查一次D1；无人进出时关门并在20t后安排E6\\
E6 & 取消过期移动事件，刷新按钮，计算目标楼层并决定上行、下行或空闲\\
E7 & 上升一层，在有请求时安排减速E2，否则继续上行\\
E8 & 下降一层，在有请求时安排减速E2，否则继续下行\\
E9 & 清除长时间停留标志并重新调用控制器\\
\bottomrule
\end{longtable}

\subsection{双电梯调度算法}

当某层出现新的方向请求时，\texttt{AssignCall}调用\texttt{DispatchElevator}。后者分别计算两部电梯的调度代价：

\[
C=10d+p+q,
\]

其中$d$为电梯与请求楼层的距离，$q$为轿厢内人数。当电梯运行方向与乘客方向不一致，或请求不在电梯当前运行前方时，附加惩罚$p=50$；空闲或顺路时$p=0$。系统选择代价最小的电梯，代价相同时选择编号较小者。

分配结果写入\texttt{assignedUp[floor]}或\texttt{assignedDown[floor]}。电梯只响应分配给自己的厅外请求，但始终响应自己的\texttt{CallCar}。当对应方向已经无人等待时，\texttt{CancelAssignment}把分配值恢复为$-1$。

\subsection{复杂度分析}

设事件表当前长度为$n$，某队列长度为$m$。事件插入、指定事件删除和指定乘客删除分别为$O(n)$、$O(n)$和$O(m)$；事件弹出和普通入队为$O(1)$。楼层数和电梯数在本实验中分别固定为5和2，因此目标楼层搜索与电梯选择可以视为常数时间。整个模拟的主要开销来自事件表的有序插入。

\section{调试分析}

\subsection{编译阶段问题}

开发过程中首先处理了头文件和类型问题。原实现同时包含队列的\texttt{common.h}与\texttt{operations.h}，而后者又包含前者，导致无头文件保护时可能重复定义。最终改为仅通过相对路径包含队列操作头文件，并为电梯头文件添加宏保护。


\subsection{运行阶段问题}

\begin{enumerate}[label=\arabic*.]
    \item \textbf{模拟时间倒退。}主循环弹出事件后没有立即更新\texttt{simulator.time}，导致控制器使用旧时间插入新事件。修正为每次分派前执行\texttt{simulator.time=time}，随后检查完整日志，时间倒退次数为0。
    \item \textbf{E6无法启动移动。}上行和下行分支曾错误嵌套在\texttt{goal==-1}内部，使有效目标无法进入移动分支。将条件重构为互斥的\texttt{if--else if--else}结构。
    \item \textbf{下行事件写成上行事件。}E8继续下降时误插入E7，导致方向突变。修正为继续插入E8。
    \item \textbf{遗漏一名离开乘客。}倒序循环条件曾写为\texttt{i>0}，遗漏数组下标0；修改为\texttt{i>=0}。
    \item \textbf{门无法关闭。}最后一个M5或M6完成后没有清除D1。增加人员事件检查，并在无人继续进出时设置$D1=0,D3=1$。
    \item \textbf{旧事件重复驱动。}空闲事件和移动事件未随状态改变取消，产生重复E1/E7/E8。通过\texttt{eventTrue}去重，并由\texttt{DeleteEvent}按电梯编号取消过期事件。
    \item \textbf{双电梯重复响应。}若只保存CallUp与CallDown，两部电梯会同时认为请求属于自己。增加\texttt{assignedUp}和\texttt{assignedDown}，保证一个方向请求只分配给一部电梯。
\end{enumerate}

\section{测试结果}

\subsection{测试结果汇总}

\begin{center}
\begin{tabular}{lcc}
\toprule
指标 & 单电梯 & 双电梯\\
\midrule
测试乘客数 & 15 & 15\\
成功到达人数 & 10 & 13\\
超时放弃人数 & 5 & 2\\
全部乘客得到最终处理 & 是 & 是\\
事件时间倒退次数 & 0 & 0\\
最后处理时刻/t & 5320 & 5209\\
结束时电梯状态 & 1层、Idle & 两部均为1层、Idle\\
\bottomrule
\end{tabular}
\end{center}

\subsection{运行结果示例}

图\ref{fig:single}和图\ref{fig:double}分别给出了单电梯和双电梯程序在15位乘客测试输入下的实际运行日志（截取自程序控制台输出）。单电梯日志显示模拟在Time=5320结束，电梯最终回到1层并处于Idle状态；双电梯日志显示模拟在Time=5209结束，两部电梯均回到1层并处于Idle状态，统计结果为Left=13、GaveUp=2、Total=15，与上表汇总数据一致。

\begin{figure}[htbp]
\centering
\includegraphics[width=0.85\textwidth]{single.png}
\caption{单电梯运行结果示例}
\label{fig:single}
\end{figure}

\begin{figure}[htbp]
\centering
\includegraphics[width=0.85\textwidth]{double.png}
\caption{双电梯运行结果示例}
\label{fig:double}
\end{figure}

双电梯运行中，0号电梯搭载8人，1号电梯搭载5人，说明调度器实际使用了两部电梯，而不是把全部请求固定分配给同一部电梯。与单电梯相比，双电梯使成功到达人数由10人增加到13人，放弃人数由5人下降到2人。对本组固定输入而言，增加电梯数量及基于方向、距离的调度降低了乘客等待超时的概率。


\section{总结}

本项目使用有序事件链表表达时间推进，使用链队列表达乘客等待与轿厢乘客集合，实现了乘客与电梯两个相互作用的状态机。双电梯把电梯状态设计成\texttt{Elevator}数组，同时保留共享楼层资源，并通过请求分配数组解决重复响应问题。


\appendix
\section{项目文件结构}

\begin{lstlisting}[language={},numbers=none]
queue/Elevator/
|-- include/
|   |-- common.h
|   |-- operations.h
|   |-- multi_common.h
|   `-- multi_operations.h
|-- src/
|   |-- operations.c
|   `-- multi_operations.c
|-- main.c
|-- multi_main.c
`-- Elevator_Report.tex

queue/single_linked_queue/
|-- include/common.h
|-- include/operations.h
`-- src/operations.c
\end{lstlisting}

\section{双电梯公共数据结构}
\lstinputlisting[caption={multi\_common.h}]{include/multi_common.h}

\section{双电梯函数声明}
\lstinputlisting[caption={multi\_operations.h}]{include/multi_operations.h}

\section{双电梯主程序}
\lstinputlisting[caption={multi\_main.c}]{multi_main.c}

\section{双电梯操作实现}
\lstinputlisting[caption={multi\_operations.c}]{src/multi_operations.c}

\section{复用的链队列接口}
\lstinputlisting[caption={single\_linked\_queue/include/operations.h}]{../single_linked_queue/include/operations.h}

\end{document}
